\section{Formal verification}\label{p2:sec-ver}
\begingroup\sloppy

The statements of this paper are verified in the Lean~4 development shared with Paper~I \cite{LeanRepo,Lean4}, with the
toolchain \texttt{v4.34.1}, Mathlib \cite{Mathlib2020} at the tag of the same name, and the unmodified Lean library
accompanying \cite{OpenAI2026}. That library contributes lemmas and a few generic differential-geometric operators
(exterior derivative, chart differential); the definitions of the Hamiltonian and of the algorithms belong to the
development, and none of the library's packing constructions or existence results is used. In the directory
\texttt{verification/}, \texttt{Problem.lean} defines the problem using Mathlib alone, \texttt{Algorithm.lean} exposes
the top level of the construction, \texttt{Statement.lean} states its correctness, and \texttt{Computability.lean} and
\texttt{EpochComputability.lean} state correctness and termination of the finite-output algorithms. Theorems in \texttt{verification/} have no
namespace; other names are relative to \texttt{SymplecticPacking}. The statement files, the trusted base and the
metaprograms are described in \PoneRef{sec:ver}.

\begin{table}[ht]
\centering
\footnotesize
\renewcommand{\arraystretch}{1.15}
\begin{tabular}{@{}L{0.24\textwidth}L{0.7\textwidth}@{}}
\toprule
Statement & Lean \\
\midrule
Definitions~\ref{def:intro-spec}, \ref{def:intro-admissible} & \texttt{Problem.PackingSpec}, \texttt{Problem.Admissible} \\
Theorem~\ref{thm:main}(i) & \texttt{main\_theorem}, \texttt{main\_time\_support}, \texttt{main\_support\_radius}; for the
 flow, \texttt{flow\_is\_integral\_curve}, \texttt{flow\_symplectic\_diffeomorphism}, \texttt{flow\_images\_disjoint} \\
Theorem~\ref{thm:main}(ii) & \texttt{\#check\_explicit General.packingHamiltonian} (a metaprogram) \\
Theorem~\ref{thm:main}(iii) & \texttt{original\_name\_algorithm}, with \texttt{\#check\_constructive} on the machine;
 finite form: \texttt{General.Compile.main\_finite} \\
Corollary~\ref{thm:main-sy} & \texttt{neighbourhood\_criterion}, \texttt{flow\_neighbourhood\_embedding} \\
Corollary~\ref{thm:main-equal} & \texttt{equal\_balls} \\
Corollary~\ref{thm:main-cone4} & \texttt{unequal\_cone} \\
Corollary~\ref{thm:main-epoch} & \texttt{epoch\_packing} \\
Theorem~\ref{thm:eff-algorithm} & \texttt{original\_name\_algorithm}, \texttt{epoch\_name\_algorithm}, \texttt{epoch\_name\_algorithm\_packing} \\
embeddings, by case & \texttt{Branch2.embeddingsEq\_spec}, \texttt{Branch2.embeddingsGen\_spec} ($n=2$, linear and scalar datum), \texttt{BranchSmall.embeddings\_spec} (all $c_i<\frac C2$),
 \texttt{BranchLargeHigh.embeddings\_spec} ($n\ge4$, one $c_j\ge\frac C2$), \texttt{BranchLarge3.embeddings\_spec}
 ($n=3$, one $c_j\ge\frac C2$) \\
transfers & \texttt{BranchSmall.vectorTransfer\_spec}, \texttt{BranchLargeHigh.shell\_pullback},
 \texttt{BranchLarge3.cubicTransfer\_spec} \\
comparison and handle & \texttt{HighDim.CmpInput.comparatorDatum}, \texttt{HighDim.liftGen\_spec},
 \texttt{HighDim.RawModelAgreement.native\_packing\_spec}, \texttt{HighDim.assembly\_spec} \\
\bottomrule
\end{tabular}
\caption{Formal counterparts of the statements of this paper. Every case of Construction~\ref{cons:ham-branch} is
covered by \texttt{main\_theorem} and the algorithm theorems.}
\label{p2:tab-ver}
\end{table}

\emph{The statement.} Phase space is $\texttt{Problem.E}~n=\texttt{EuclideanSpace}\ \R\ (\texttt{Fin}\ (2n))$;
\texttt{omega0} and \texttt{hamVF} give $\omega_0$ and $X_H=(\partial_yH,-\partial_xH)$, so $\iota_{X_H}\omega_0=dH$.
\texttt{capBall}~$s$~$c=\{z:\pi\|z-s\|^2\le c\}$ and \texttt{capBallOpen} is its strict version. \texttt{Admissible},
\texttt{BallsDisjoint} and \texttt{PackingSpec} are Definitions~\ref{def:intro-admissible} and~\ref{def:intro-spec} and the
disjointness hypothesis, field by field; \texttt{PackingSpec} quantifies over every integral curve on $[0,1]$, so it
does not depend on a chosen solution. \texttt{Algorithm.lean} shows by \texttt{rfl} that
\texttt{General.packingHamiltonian} is \texttt{HamEnd.hamiltonianOfEmbeddings} applied to
\texttt{General.selectEmbeddings}, the case distinction of Construction~\ref{cons:ham-branch}. The Lean statements agree
with the text, with these differences of form:
\texttt{equal\_balls} does not assume $\eps>0$; the algorithm theorems assert \texttt{Nat.RecursiveIn} relative to the
oracles given by the input names, and that one program serves every input is certified by
\texttt{\#check\_constructive}; \texttt{flow\_neighbourhood\_embedding} uses the neighbourhood $(\phi^1_H)^{-1}(\Bo(C))$
and holds for every $n\ge2$.

\emph{What is verified.} The kernel checks the theorems of \texttt{Statement.lean} and \texttt{Computability.lean}
and all their dependencies in the development, the vendored library and Mathlib. The trusted base is the kernel, the
definitions in \texttt{Problem.lean} and the axioms \texttt{propext}, \texttt{Classical.choice} and
\texttt{Quot.sound}; the build fails if the transitive axiom set of any of these theorems contains anything else, in
particular \texttt{sorryAx}. Explicitness (level L1) and uniformity (level L3) are checked by the metaprograms
\texttt{\#check\_explicit} and \texttt{\#check\_constructive}, which are not themselves verified. The first is described
in \PoneRef{sec:ver} and \PoneRef{sec:xpl}; it is also applied to the algorithm terms. The second, applied to the two
oracle machines, additionally forbids any classical constant, including classical decidability, so each is one fixed
program in its oracle streams and cannot branch on an undecidable property of the named reals.
\texttt{test/ExplicitTest.lean} contains terms both commands must reject, among them a machine branching on
$x<\frac13$ for real $x$.

\emph{Independence.} The command \texttt{\#assert\_construction\_independence} in
\texttt{verification/Audit/Independence.lean} traverses types, definitions and proof bodies of the packing, flow,
corollary and algorithm theorems and fails if they reach the vendored packing predicate or existence theorems, the
holomorphic-line and two-ball obstruction proofs, or the curve-analysis modules of the vendored library. In
particular the sufficiency direction of Corollary~\ref{thm:main-sy} does not use the packing existence theorem of
\cite{OpenAI2026}. The converse is checked separately: it may use the classical obstructions, but no sufficiency
theorem.

\emph{Source tags.} The sources of the main text and the supplement attach Lean declaration names to further
statements; \texttt{verification/PAPER-INDEX.md} indexes them, and the build checks that every tagged name exists and
depends only on the three axioms. A tag is a reference to related declarations, without a claim that the Lean
statement covers the printed one, except as follows. Definitions~\ref{def:intro-spec} and~\ref{def:intro-admissible},
Theorem~\ref{thm:main}(i), (iii) and Corollaries~\ref{thm:main-sy}--\ref{thm:main-epoch} are verified in full, up to
the differences of form above; Theorem~\ref{thm:main}(ii) is checked by \texttt{\#check\_explicit}.
For Theorem~\ref{thm:cmp-realisation} (\Sref{Theorem}{thm:plc-lift}), \Sref{Theorem}{I:thm:pln-isotopy} and
\Sref{Lemma}{I:lem:rtc-ball-chart} the formal statement is the consequence used by the construction, not every
assertion or the whole domain of the printed statement; the restrictions are recorded in
\PoneRef{sec:xpl-interface}. Table~\ref{tab:guide-lean} lists the Lean declarations for the checking route of
\Ssec{p2:sec-guide}.

{\footnotesize
\renewcommand{\arraystretch}{1.12}
\begin{longtable}{@{}p{0.08\textwidth}>{\raggedright\arraybackslash}p{0.86\textwidth}@{}}
\caption{Lean declarations for the steps H1--H7 of \Ssec{p2:sec-guide} and the steps F1, F9--F11 of Paper~I (\PoneRef{sec:guide-four}). Names are relative to \texttt{SymplecticPacking}, except the theorems of
\texttt{verification/}, which have no namespace.}\label{tab:guide-lean}\\
\toprule
step & Lean declarations\\\midrule
\endfirsthead
\multicolumn{2}{c}{\tablename\ \thetable\ (continued)}\\
\toprule
step & Lean declarations\\\midrule
\endhead
\midrule
\multicolumn{2}{r}{Continued on the next page}\\
\endfoot
\bottomrule
\endlastfoot
F1 & \texttt{General.selectEmbeddings}, \texttt{Branch2.EqualCase}\\
F9 & \texttt{General.Ham.Isotopy.IsoInput.commonIsotopy}, \texttt{\dots isoF\_sympl}, \texttt{\dots isoF\_disj}\\
F10 & \texttt{HamEnd.hamiltonianOfEmbeddings}, \texttt{General.Ham.hamVF\_ofIsotopy\_image}, \texttt{HamEnd.hamiltonian\_end}\\
F11 & \texttt{main\_theorem}, \texttt{four\_dim\_theorem}, \texttt{epoch\_packing}, \texttt{epoch\_theorem} in
 \texttt{verification/Statement.lean}\\
H1 & \texttt{HighDim.CmpInput.comparatorDatum}, \texttt{HighDim.ComparisonDatum}, \texttt{HighDim.PlanarGen.Input}\\
H2 & \texttt{Branch2.cmpInput}, \texttt{Branch2.cubicDataOf}, \texttt{Branch2.Hyp.cubicDataOf\_valid}, \texttt{Branch2.embeddingsGen\_spec}\\
H3 & the tags of \Sref{Theorem}{thm:shd-packing} in \texttt{PAPER-INDEX.md}; composition \texttt{HighDim.assembly\_spec}\\
H4 & \texttt{BranchSmall.vectorTransfer\_spec}, \texttt{BranchSmall.embeddings\_spec}\\
H5 & \texttt{BranchLargeHigh.shell\_pullback}, \texttt{BranchLargeHigh.embeddings\_spec}\\
H6 & \texttt{BranchLarge3.cubicTransfer\_spec}, \texttt{BranchLarge3.embeddings\_spec}; corollaries
 \texttt{main\_theorem}, \texttt{neighbourhood\_criterion}, \texttt{equal\_balls}\\
H7 & \texttt{General.Compile.main\_finite}, \texttt{original\_name\_algorithm}, \texttt{epoch\_name\_algorithm}\\
\end{longtable}
}
\endgroup
